% Ad Atticum 16.2 (ad-atticum-16-02)
% Source: https://raw.githubusercontent.com/PerseusDL/canonical-latinLit/master/data/phi0474/phi057/phi0474.phi057.perseus-lat1.xml
% Fetched by scripts/fetch_latin.py

% Dateline (Perseus): Scr. in Puteolano v Id. Quint. a. 710 (44).

\ciceroLetterOpener{CICERO ATTICO salutem}

\ciceroSection{1} vi Idus duas epistulas accepi, unam a meo tabellario, alteram a Bruti. de Buthrotiis longe alia fama in his locis fuerat, sed cum aliis multis hoc ferendum. Erotem remisi citius quam constitueram, ut esset qui Hortensio et †quia et quibus quidem ait se Idibus constituisse. Hortensius vero impudenter. nihil enim debetur ei nisi ex tertia pensione quae est Kal. Sext.; ex qua pensione ipsa maior pars est ei soluta aliquanto ante diem. sed haec Eros videbit Idibus. de Publilio autem, quod perscribi oportet, moram non puto esse faciendam. sed cum videas quantum de iure nostro decesserimus qui de residuis C_C_C_C_ HS C_C_ praesentia solverimus, reliqua rescribamus, loqui cum eo, si tibi videbitur, poteris eum commodum nostrum exspectare debere, cum tanta sit a nobis iactura facta iuris.

\ciceroSection{2} sed amabo te, mi Attice, (videsne quam blande?), omnia nostra, quoad eris Romae, ita gerito, regito, gubernato ut nihil a me exspectes. quamquam enim reliqua satis apta sunt ad solvendum, tamen fit saepe ut ii qui debent non respondeant ad tempus. si quid eius modi acciderit, ne quid tibi sit fama mea potius. non modo versura verum etiam venditione, si ita res coget, nos vindicabis.

\ciceroSection{3} Bruto tuae litterae gratae erant. fui enim apud illum multas horas in Neside, cum paulo ante tuas litteras accepissem. delectari mihi Tereo videbatur et habere maiorem Accio quam Antonio gratiam. mihi autem quo laetiora sunt, eo plus stomachi et molestiae est populum Romanum manus suas non in defendenda re publica sed in plaudendo consumere. mihi quidem videntur istorum animi incendi etiam ad repraesentandam improbitatem suam. sed tamen dúm modo doleant áliquid, doleant quídlibet.

\ciceroSection{4} consilium meum quod ais cotidie magis laudari non moleste fero exspectabamque si quid de eo ad me scriberes. ego enim in varios sermones incidebam. quin etiam idcirco trahebam ut quam diutissime integrum esset. sed quoniam furcilla extrudimur, Brundisium cogito. facilior enim et exploratior devitatio legionum fore videtur quam piratarum qui apparere dicuntur. Sestius vi Idus exspectabatur sed non venerat, quod sciam. Cassius cum classicula sua venerat. ego cum eum vidissem, v Id. in Pompeianum cogitabam, inde Aeculanum. nosti reliqua. de Tutia ita putaram.

\ciceroSection{5} de Aebutio non credo nec tamen curo plus quam tu. Planco et Oppio scripsi equidem quoniam rogaras, sed, si tibi videbitur, ne necesse habueris reddere. cum enim tua causa fecerint omnia, vereor ne meas litteras supervacaneas arbitrentur. Oppio quidem utique quem tibi amicissimum cognovi. verum ut voles.

\ciceroSection{6} tu quoniam scribis hiematurum te in Epiro, feceris mihi gratum si ante eo veneris quam mihi in Italiam te auctore veniendum est. litteras ad me quam saepissime; si de rebus minus necessariis, aliquem nactus; sin autem erit quid maius, domo mittito. Ἡρακλείδειον , si Brundisium salvi, adoriemur. de gloria misi tibi. custodies igitur, ut soles, sed notentur ἐκλογαὶ quas Salvius bonos auditores nactus in convivio dumtaxat legat. mihi valde placent, mallem tibi. etiam atque etiam vale.
